\section{Remarks}
Overall, there is no single solution to the issues initially outlined. When multiple fragility or vulnerability models are defined using different \acp{IM}, analysts must choose an approach to enable comparison and collective application; in such cases, methods based on total probability should be preferred. Historically, the option of adopting a single \ac{IM} has been common, with $PGA$ being the most widely used over recent decades. However, efficiency and sufficiency considerations have highlighted significant limitations in using $PGA$ to characterise seismic vulnerability. The results presented here further demonstrate how restrictive model definition becomes when using this \ac{IM}, since it complicates fragility function conversion and can substantially overestimate losses when converted models are applied. An alternative is to use next-generation \acp{IM} such as $Sa_{avg}(T)$ or $FIV3(T)$, which retain the practicality of adopting a unique \ac{IM} while avoiding many of the biases and constraints associated with $PGA$. This reinforces the case for future fragility and vulnerability modelling to incorporate these improved \acp{IM}.

The conversion methods investigated employed different strategies, ranging from assumptions of uniform hazard to uniform risk and also the application of total probability. The total probability approach performed best, although it requires more detailed seismic hazard information, which is often available but not always, confirming earlier observations by \citet{Suzuki2020}. Methods 2a, 2b, and 2c, however, indicated that a comprehensive description of hazard disaggregation information is not always needed. Simpler methods can work at times (i.e., Methods 3, 4 and 5), but their success depends heavily on the relative efficiency of the original and target \acp{IM}, which governs the reliability of the conversion. Methods 3 and 4 were expected to perform modestly, which was confirmed. Greater expectations were placed on the feasibility of Method 5, as it is grounded in the principle of uniform risk equivalence; however, this ultimately did not materialise into a usable method. Table \ref{tab:method_comparison} gives an overall summary of each method, noting their required hazard information and other additional models/inputs, their perceived implementation complexity in addition to the relative accuracy observed for the results presented here, provided the initial $IM_1$ is an efficient one. Inevitably, these methods require computational tools for implementation, and some of these tools are provided in the associated GitHub repository, although the challenge of characterising suitable hazard inputs would nevertheless remain to be addressed. 
